\chapter{設計と実装}
\label{chap:design}

本章では，第\ref{chap:proposal}章で示した提案アーキテクチャを実現するための詳細設計，プロトタイプ実装，およびプロトコル仕様上の制約を述べる．

第\ref{chap:proposal}章では，4G/5Gの機能共通性分析に基づき論理アーキテクチャを導出し，主要な設計課題（手順タイミングの不一致，認証・鍵管理の分離，QoSモデルの差異）を特定した．本章では，これらの課題に対する具体的な設計解を示し，概念実証（PoC）として実際に実装した範囲・簡略化・論理設計との差分を明らかにし，プロトコル仕様上の制約を整理する．具体的には，第\ref{sec:detailed_design}節で詳細設計を，第\ref{sec:prototype_implementation}節でプロトタイプ実装を，第\ref{sec:protocol_constraints}節で制約事項を述べる．

\section{詳細設計}
\label{sec:detailed_design}

\subsection{C-Plane変換ロジック}

\subsubsection{S1 Setupシーケンス（ID変換）}
eNBの起動時に行われるS1 Setup手順において，ゲートウェイは以下のID変換を行う．
\begin{itemize}
  \item \textbf{Global eNB ID $\rightarrow$ Global gNB ID}: 4Gの基地局IDを5Gの基地局ID形式にマッピングする．S1SetupRequestからPLMN（MCC/MNC）およびeNB-IDを抽出し，対応する形式のgNB-IDを生成する．
  \item \textbf{TAC}: S1SetupRequestのSupportedTAsからTACを抽出して設定し，エリア情報を維持する．
\end{itemize}
これにより，AMFは接続元を「5G基地局（gNB）」として認識し，NG Setupを完了させる．

\subsubsection{Registrationシーケンス（NASメッセージ変換）}
UEからのAttach Request（4G NAS）は，ゲートウェイによって5G NASメッセージ（Registration Request）に変換され，AMFへ転送される．
\begin{itemize}
	\item \textbf{Initial UE Message}: S1APのInitial UE MessageをNGAPのInitial UE Messageに変換．
	\item \textbf{NAS Transport}: 4G NASコンテナ内の情報を抽出し，5G NAS形式に変換する．変換不可能なメッセージについては，ログ出力の上で当該処理を中止する．
	\item \textbf{認証・セキュリティ}: ゲートウェイが事前に保持するUEの秘密鍵（Ki\footnote{SIMに格納される加入者固有の共有秘密鍵．}, OPc\footnote{SIMカードの製造者鍵OPとKiから派生したオペレータ固有値．}）を用いて，AMFからの認証要求に対する応答（RES*生成）を代行し，セキュリティコンテキストを確立する．
\end{itemize}

\subsubsection{NASメッセージ変換対応表}
表\ref{tab:nas_conversion_ul}および表\ref{tab:nas_conversion_dl}に，主要なNASメッセージの変換対応を示す．

\begin{table}[htbp]
\centering
\caption{NASメッセージ変換対応表（上り：UE→CN）}
\label{tab:nas_conversion_ul}
\begin{tabular}{|l|l|}
\hline
\textbf{4G NAS} & \textbf{5G NAS} \\
\hline
Attach Request & Registration Request \\
\hline
Auth Response (RES) & Auth Response (RES*へ変換) \\
\hline
Security Mode Complete & Security Mode Complete \\
\hline
Attach Complete + ESM情報 & Registration Complete + PDU Session Est Req \\
\hline
\end{tabular}
\end{table}

\begin{table}[htbp]
\centering
\caption{NASメッセージ変換対応表（下り：CN→UE）}
\label{tab:nas_conversion_dl}
\begin{tabular}{|l|l|}
\hline
\textbf{5G NAS} & \textbf{4G NAS} \\
\hline
Auth Request (ngKSI) & Auth Request (eKSI) \\
\hline
Security Mode Cmd (NEA/NIA) & Security Mode Cmd (EEA/EIA) \\
\hline
Registration Accept + 5G-GUTI & Attach Accept + GUTI \\
\hline
PDU Session Est Accept (5QI, DNN) & Act Def EPS Bearer (QCI, APN) \\
\hline
\end{tabular}
\end{table}

また，ID変換については表\ref{tab:id_conversion}に示す対応関係に基づいて行う．

\begin{table}[htbp]
\centering
\caption{ID変換対応表}
\label{tab:id_conversion}
\begin{tabular}{|l|l|}
\hline
\textbf{4G} & \textbf{5G} \\
\hline
IMSI & SUCI (NULL scheme) \\
\hline
GUTI & 5G-GUTI \\
\hline
eKSI (3bit) & ngKSI (4bit, TSC=0) \\
\hline
\end{tabular}
\end{table}

\subsubsection{PDU Session Establishment（Early S1 Setup）}
\label{sec:pdu_session_establishment}
本提案の核心となる独自シーケンスである．

\paragraph{課題}
5GではRegistration完了後にPDUセッション確立が行われるのに対し，4GではAttach手順の一部としてDefault Bearer（データ経路）が確立される．4G UEはAttach Acceptを受信した時点でデータ通信可能な状態を期待するため，5G側のPDUセッション確立タイミングとの間にズレが生じる．このタイミングのズレを解消するため，以下の「Early S1 Setup」手順を導入した．

\begin{enumerate}
	\item \textbf{Registration Accept受信時}: AMFからのRegistration Acceptを受信した時点で，ゲートウェイはUEに対して「Attach Accept」と共に「Initial Context Setup Request (S1)」を先行して送信する．
	\item \textbf{仮TEIDの割り当て}: この時点ではUPF側のTEIDが未定であるため，ゲートウェイは自身のTEID（S1N2 TEID）をeNBに通知し，U-Planeの経路を仮確立する．
	\item \textbf{PDU Session Establishment Request}: Registration完了後，ゲートウェイはUEに代わってPDUセッション確立要求を生成し，AMFへ送信する．UEから受信したAttach RequestのESM情報（APN等）を解析して使用する．
	\item \textbf{TEID Update}: UPFから正規のTEIDが通知された後，ゲートウェイは内部のマッピングテーブルを更新し，eNBからのパケットをUPFへ転送可能にする．
\end{enumerate}

図\ref{fig:sequence_diagram}に，実際のプロトタイプ動作におけるシーケンス図を示す．ここでは，Registration Accept受信時点（Registration Complete送信前）にEarly S1 Setupが実行され，その後Registration Complete検出を契機にゲートウェイがPDU Session Establishment RequestをAMFへ送信し，AMFからNGAP InitialContextSetupRequest（PDU Session Resource Setupを含む）が返される様子を示している．

\begin{figure}[htbp]
  \centering
  \begin{tikzpicture}[scale=0.85, transform shape, every node/.style={font=\small}]
    % Nodes
    \node (enb) at (0,0) [draw, rounded corners, fill=GoogleBlue!20, minimum width=2cm] {eNB};
    \node (conv) at (6,0) [draw, rounded corners, fill=GoogleYellow!20, minimum width=2cm] {Converter};
    \node (amf) at (12,0) [draw, rounded corners, fill=GoogleGreen!20, minimum width=2cm] {AMF};

    % Lifelines
    \draw [thick, dotted] (enb) -- (0,-14.5);
    \draw [thick, dotted] (conv) -- (6,-14.5);
    \draw [thick, dotted] (amf) -- (12,-14.5);

    % Time steps
    \footnotesize

    % Step 1: Attach / Reg
    \draw [->, thick] (0,-1) -- (6,-1) node [midway, above] {Initial UE Msg (Attach Req)};
    \draw [->, thick] (6,-1.5) -- (12,-1.5) node [midway, above] {Initial UE Msg (Reg Req)};

    % Step 2: Auth
    \draw [<-, thick] (6,-2.5) -- (12,-2.5) node [midway, above] {Auth Request};
    \draw [<-, thick] (0,-2.5) -- (6,-2.5) node [midway, above] {Auth Request};

    \draw [->, thick] (0,-3.5) -- (6,-3.5) node [midway, above] {Auth Response};
    \draw [->, thick] (6,-3.5) -- (12,-3.5) node [midway, above] {Auth Response};

    % Step 3: Security
    \draw [<-, thick] (6,-4.5) -- (12,-4.5) node [midway, above] {Sec Mode Command};
    \draw [<-, thick] (0,-4.5) -- (6,-4.5) node [midway, above] {Sec Mode Command};

    \draw [->, thick] (0,-5.5) -- (6,-5.5) node [midway, above] {Sec Mode Complete};
    \draw [->, thick] (6,-5.5) -- (12,-5.5) node [midway, above] {Sec Mode Complete + Reg Req};

    % Step 4: Registration Accept from AMF
    \draw [<-, thick] (6,-6.5) -- (12,-6.5) node [midway, above] {DL NAS (Reg Accept)};

    % Step 5: Early S1 Setup (Converter initiates upon Reg Accept detection)
    \node [draw, fill=GoogleYellow!30, align=center, font=\scriptsize] at (-1.5,-7.9) {Early S1 Setup\\(Reg Accept検出時)};
    \draw [<-, thick] (0,-7.5) -- (6,-7.5) node [midway, above] {Init Ctxt Setup Req + Attach Accept};
    \draw [->, thick] (0,-8.3) -- (6,-8.3) node [midway, above] {Init Ctxt Setup Resp};

    % Step 6: Attach Complete / Reg Complete
    \draw [->, thick] (0,-9.3) -- (6,-9.3) node [midway, above] {UL NAS (Attach Cmp)};
    \draw [->, thick] (6,-9.8) -- (12,-9.8) node [midway, above] {UL NAS (Reg Cmp)};

    % Step 7: PDU Session Establishment (Converter generates on behalf of UE)
    \node [draw, fill=GoogleYellow!30, align=center, font=\scriptsize] at (13.5,-10.8) {Reg Cmp検出後\\PDU Session要求};
    \draw [->, thick] (6,-11.0) -- (12,-11.0) node [midway, above] {PDU Session Est Req};

    % Step 8: AMF responds with PDU Session setup
    \draw [<-, thick] (6,-12.2) -- (12,-12.2) node [midway, above] {NGAP ICS Req (PDU Session)};
    \draw [->, thick] (6,-13.2) -- (12,-13.2) node [midway, above] {NGAP ICS Resp};

    % Time axis label
    \node at (-1, -7) [rotate=90] {Time};
    \draw [->] (-0.5, -1) -- (-0.5, -14);

    % 凡例（下部中央）
    \node[draw, fill=white, anchor=north, font=\scriptsize, inner sep=4pt] at ([yshift=-8pt]current bounding box.south) {
      \begin{tabular}{@{}l@{\hspace{1em}}l@{\hspace{1em}}l@{\hspace{1em}}l@{}}
      \csname tikz\endcsname[baseline=-0.6ex]\fill[GoogleBlue!20] (0,0) rectangle (0.18,0.18);\ eNB &
      \csname tikz\endcsname[baseline=-0.6ex]\fill[GoogleYellow!20] (0,0) rectangle (0.18,0.18);\ 変換器 &
      \csname tikz\endcsname[baseline=-0.6ex]\fill[GoogleGreen!20] (0,0) rectangle (0.18,0.18);\ AMF &
      \csname tikz\endcsname[baseline=-0.6ex]\fill[GoogleYellow!30] (0,0) rectangle (0.18,0.18);\ 能動処理
      \end{tabular}
    };

  \end{tikzpicture}
  \caption{プロトタイプにおける接続シーケンス（Early S1 Setupを含む）}
  \label{fig:sequence_diagram}
\end{figure}

\subsection{U-Plane変換とTEIDマッピング}
U-Planeでは，eNBとUPF間のユーザデータ転送を仲介する．4GのS1-Uインタフェースと5GのN3インタフェースは共にGTPv1-Uプロトコルを使用するため，プロトコル変換は不要であるが，TEIDと宛先IPアドレスの書き換えが必要となる．

なお，実装によっては方向別（Uplink用／Downlink用）のマッピングテーブルを持ち，各テーブル内でTEIDが一意であることを前提として，GTP-Uヘッダ内のTEIDのみを検索キーとしてマッピングを行う方式も有効である．

\subsubsection{TEIDマッピングテーブル}
ゲートウェイは以下のTEID変換テーブルを動的に管理する．

\begin{table}[h]
\centering
\caption{TEIDマッピングテーブル}
\label{tab:teid_mapping}
\small
\begin{tabularx}{\textwidth}{|l|X|X|}
\hline
方向 & 検索キー (TEID) & 出力 (Dst IP, TEID) \\
\hline
Uplink (eNB $\rightarrow$ 5GC) & S1N2 TEID \newline （S1-U側：GW割当） & (UPF IP, N3 TEID) \\
\hline
Downlink (5GC $\rightarrow$ eNB) & DLパケットのTEID \newline （eNB downlink TEIDと同値） & (eNB IP, TEID) \newline ※TEID書換なし \\
\hline
\end{tabularx}
\end{table}

パケット受信時，ゲートウェイはGTP-Uヘッダ内のTEIDを参照して転送処理を行う．UplinkではTEIDマッピングテーブルを参照し，GTP-UヘッダのTEIDフィールドをN3 TEIDに書き換えてUPFへ送出する．Downlinkでは，UPFから受信したパケットのTEID（ゲートウェイがUPFに通知したTEIDであり，eNBのdownlink TEIDと同値に設定）を用いてUEコンテキストを特定し，対応するeNBのS1-U IPを取得して転送する．転送はOSのネットワークスタックを介して行われるため，外部IP/UDPヘッダは新規に生成され，チェックサムもOSにより計算される．

\subsubsection{マッピングエントリの生成}
TEIDマッピングエントリは，C-Planeのセッション確立手順と連動して以下のタイミングで生成される．

\begin{enumerate}
\item \textbf{S1N2 TEIDの事前割当}: \ref{sec:pdu_session_establishment}節の「Early S1 Setup」手順において，ゲートウェイが管理するTEID（S1N2 TEID）をeNBに通知する．
\item \textbf{eNB側TEIDの取得}: eNBからのInitial Context Setup Responseに含まれるS1-U TEIDを記録する．
\item \textbf{N3側TEIDの取得}: PDU Session Resource Setup手順により通知されるUPF側のN3 TEIDを取得し，Uplink用のマッピングエントリを完成させる．
\end{enumerate}

この設計により，eNBとUPFはそれぞれ対向ノードとしてゲートウェイのみを認識し，互いの存在（IPアドレス体系の違いなど）を意識することなく通信が可能となる．

\subsubsection{GTP-Uパケット処理}
具体的なパケット処理フローを以下に示す．

\paragraph{Uplink処理}
\begin{enumerate}
\item eNBからGTP-Uパケットを受信（宛先TEID = S1N2 TEID）
\item TEIDマッピングテーブルを検索し，対応するUPF情報を取得
\item GTP-UヘッダのTEIDをN3 TEID（上り：UPF側）に書き換え
\item UPFをUDP送信先として送出（外側IP/UDPヘッダはOSスタックで生成）
\end{enumerate}

\paragraph{Downlink処理}
Downlinkでは，受信パケットのTEIDをキーにUEコンテキストを参照し，対応するeNBへ転送する．
\begin{enumerate}
\item UPFからGTP-Uパケットを受信（TEID = eNBのdownlink TEID）
\item UEコンテキストを参照し，送信先eNB（S1-U IP）を決定
\item eNBをUDP送信先として送出（外側IP/UDPヘッダはOSスタックで生成）
\end{enumerate}

図\ref{fig:uplane_processing}に，U-Planeにおけるパケット処理とヘッダ変換の概要を示す．Uplinkでは変換器がeNBから受信したパケットのTEIDをN3用に書き換える．DownlinkではUPFがeNBのS1-U TEIDで送出するため，変換器はTEIDを書き換えず宛先IPのみeNBに変更する．

\begin{figure}[htbp]
  \centering
  \begin{tikzpicture}[scale=0.85, transform shape, every node/.style={font=\small}]
    % Uplink Flow
    \node[font=\bfseries] at (0, 2.5) {Uplink (eNB $\rightarrow$ UPF)};

    % Packet 1 (eNB -> GW)
    \node[draw, fill=GoogleBlue!10, minimum width=3.5cm, minimum height=0.8cm, align=center] (pkt1) at (0, 1.5) {IP: GW / UDP: 2152\\GTP: \textbf{S1-U TEID (UL)}};
    \node[left] at (-2, 1.5) {From eNB};

    % Processing Box
    \node[draw, thick, fill=GoogleYellow!10, minimum width=3cm, minimum height=1.2cm, align=center] (proc_ul) at (5, 1.5) {TEID Mapping\\Search \& Rewrite};
    \draw[->, thick] (pkt1) -- (proc_ul);

    % Packet 2 (GW -> UPF)
    \node[draw, fill=GoogleGreen!10, minimum width=3.5cm, minimum height=0.8cm, align=center] (pkt2) at (10, 1.5) {IP: UPF / UDP: 2152\\GTP: \textbf{N3 TEID}};
    \node[right] at (12, 1.5) {To UPF};
    \draw[->, thick] (proc_ul) -- (pkt2);

    % Downlink Flow
    \node[font=\bfseries] at (0, -1.5) {Downlink (UPF $\rightarrow$ eNB)};

    % Packet 3 (UPF -> GW)
    \node[draw, fill=GoogleGreen!10, minimum width=3.5cm, minimum height=0.8cm, align=center] (pkt3) at (0, -2.5) {IP: GW / UDP: 2152\\GTP: \textbf{eNB S1-U TEID}};
    \node[left] at (-2, -2.5) {From UPF};

    % Processing Box
    \node[draw, thick, fill=GoogleYellow!10, minimum width=3cm, minimum height=1.2cm, align=center] (proc_dl) at (5, -2.5) {Transparent Proxy\\(Dst IP Change)};
    \draw[->, thick] (pkt3) -- (proc_dl);

    % Packet 4 (GW -> eNB)
    \node[draw, fill=GoogleBlue!10, minimum width=3.5cm, minimum height=0.8cm, align=center] (pkt4) at (10, -2.5) {IP: eNB / UDP: 2152\\GTP: \textbf{eNB S1-U TEID}};
    \node[right] at (12, -2.5) {To eNB};
    \draw[->, thick] (proc_dl) -- (pkt4);

    % 凡例（下部中央）
    \node[draw, fill=white, anchor=north, font=\scriptsize, inner sep=4pt] at ([yshift=-8pt]current bounding box.south) {
      \begin{tabular}{@{}l@{\hspace{1.5em}}l@{\hspace{1.5em}}l@{}}
      \csname tikz\endcsname[baseline=-0.6ex]\fill[GoogleBlue!10] (0,0) rectangle (0.18,0.18);\ eNB側パケット &
      \csname tikz\endcsname[baseline=-0.6ex]\fill[GoogleGreen!10] (0,0) rectangle (0.18,0.18);\ UPF側パケット &
      \csname tikz\endcsname[baseline=-0.6ex]\fill[GoogleYellow!10] (0,0) rectangle (0.18,0.18);\ 変換器内処理
      \end{tabular}
    };

  \end{tikzpicture}
  \caption{U-Planeパケット処理とヘッダ変換フロー}
  \label{fig:uplane_processing}
\end{figure}

\subsection{認証と鍵管理}
\label{subsec:auth_key_management}

\subsubsection{認証フローと鍵階層}
5G AKAでは，以下の鍵階層に従って暗号鍵が導出される．

\begin{enumerate}
\item \textbf{ルート鍵}: USIMに保存されたK（共有秘密鍵）
\item \textbf{中間鍵}: Milenage\footnote{3GPPで標準化された認証・鍵生成アルゴリズム群（TS 35.206）．}アルゴリズムによりf3(K, RAND\footnote{認証センターが生成する乱数．認証ベクトルの一部．}), f4(K, RAND)からCK, IK\footnote{Cipher Key / Integrity Key: 認証時に生成される暗号化鍵と完全性保護鍵．}を生成
\item \textbf{KAUSF}: KDF\footnote{Key Derivation Function: 入力鍵から新たな鍵を導出する関数．}(CK||IK, SN name, SQN\footnote{Sequence Number: 認証ベクトルの再利用攻撃を防ぐためのシーケンス番号．}⊕AK)により導出
\item \textbf{KSEAF, KAMF}: 順次KDF適用により生成
\item \textbf{NAS鍵}: KNASenc（暗号化用），KNASint（完全性保護用）
\end{enumerate}

5G AKAでは，CK/IKはHSS/UDMとUSIM間でのみ生成され，ネットワーク上には流れない．通常，4G-5G間のハンドオーバーではMMEとAMFがN26インターフェースを介して鍵コンテキストを共有するが，本提案のような単独のプロトコル変換器構成ではその手段がない．
そのため，変換器がUEの秘密情報（Ki, OPc）を何らかの形で保持し，擬似的な認証センターとして振る舞う設計が必要となる．
変換器は，AMFから送られてくる認証リクエスト（RAND, AUTN\footnote{Authentication Token: 認証ベクトルの一部で，SQNとMACを含む．}）を受信し，保持しているKi, OPcを用いてMilenageアルゴリズムを実行することで，CK, IKを独自に導出する．

その後，導出したCK, IKを用いて以下の式により4G互換の\textbf{KASME}を生成する（TS 33.401\cite{3GPP_TS33401} Annex A.2）．

\texttt{KASME = KDF(CK||IK, FC=0x10, SN\_id, SQN⊕AK)}

ここで，SQN⊕AKは5G認証リクエストのAUTNフィールドから抽出される．この処理により，変換器は4G側のeNBとの通信に必要な鍵コンテキストを確立できる．

図\ref{fig:auth_key_derivation}に，本提案における認証フローと鍵導出のプロセスを示す．

\begin{figure}[htbp]
  \centering
  \begin{tikzpicture}[scale=0.9, transform shape, every node/.style={font=\small}]
    % Entities
    \node[draw, fill=GoogleGreen!20, minimum width=2cm] (amf) at (0, 7) {AMF};
    \node[draw, fill=GoogleYellow!20, minimum width=2.5cm, minimum height=1.5cm, align=center] (gw) at (6, 7) {Converter\\(Local Ki/OPc)};
    \node[draw, fill=GoogleBlue!20, minimum width=2cm] (enb) at (12, 7) {eNB / UE};

    % Lifelines
    \draw[thick, dotted] (amf) -- (0, 2.5);
    \draw[thick, dotted] (gw) -- (6, 2.5);
    \draw[thick, dotted] (enb) -- (12, 2.5);

    % Step 1: Auth Req from AMF to GW
    \draw[->, thick] (0, 5.8) -- (6, 5.8) node[midway, above] {Auth Req (RAND, AUTN)};

    % Step 2: Auth Req to eNB (4G)
    \draw[->, thick] (6, 5.3) -- (12, 5.3) node[midway, above] {Auth Req (RAND, AUTN)};
    \node[below, font=\scriptsize] at (9, 5.1) {(Transferred as 4G NAS)};

    % Step 3: Auth Resp from UE
    \draw[<-, thick] (6, 4.0) -- (12, 4.0) node[midway, above] {Auth Resp (RES)};

    % Step 4: Auth Resp to AMF
    \draw[<-, thick] (0, 3.5) -- (6, 3.5) node[midway, above] {Auth Resp (RES*)};
    \node[below, font=\scriptsize] at (3, 3.3) {(RES* derived locally)};

    % Internal Process in GW (positioned to the right, outside message flow)
    \node[draw, dashed, fill=white, align=left, font=\scriptsize] (calc) at (10.5, 0.5) {
      \textbf{Internal Calculation (at Converter):}\\
      1. Extract RAND, AUTN\\
      2. Load Ki, OPc (Local)\\
      3. Run Milenage(Ki, RAND)\\
      \hspace{1em}$\rightarrow$ CK, IK, RES, AK\\
      4. Extract SQN$\oplus$AK from AUTN\\
      5. \textbf{Derive RES*} (5G KDF)\\
      \hspace{1em}$= \text{KDF}(CK||IK, \text{RES}, \text{SN\_name})$\\
      6. \textbf{Derive K\_ASME} (4G KDF)\\
      \hspace{1em}$= \text{KDF}(CK||IK, \text{SN\_id}, \text{SQN}\oplus\text{AK})$
    };
    \draw[->, dashed] (6, 4.8) -- (6, 4.3) -- (calc.west);

    % 凡例（下部中央）
    \node[draw, fill=white, anchor=north, font=\scriptsize, inner sep=4pt] at ([yshift=-8pt]current bounding box.south) {
      \begin{tabular}{@{}l@{\hspace{1em}}l@{\hspace{1em}}l@{\hspace{1em}}l@{}}
      \csname tikz\endcsname[baseline=-0.6ex]\fill[GoogleGreen!20] (0,0) rectangle (0.18,0.18);\ AMF &
      \csname tikz\endcsname[baseline=-0.6ex]\fill[GoogleYellow!20] (0,0) rectangle (0.18,0.18);\ 変換器 &
      \csname tikz\endcsname[baseline=-0.6ex]\fill[GoogleBlue!20] (0,0) rectangle (0.18,0.18);\ eNB/UE &
      \csname tikz\endcsname[baseline=-0.6ex]\draw[thick, dashed] (0,0) rectangle (0.18,0.18);\ 計算処理
      \end{tabular}
    };

  \end{tikzpicture}
  \caption{認証フローと鍵導出プロセス（ゲートウェイによるK\_ASME生成）}
  \label{fig:auth_key_derivation}
\end{figure}

\subsubsection{SQN同期問題への対処}
5G-to-4G変換における重要な課題は，シーケンス番号（SQN）の同期である．UEとネットワーク側のSQNが不一致の場合，認証失敗（Authentication Failure with Synch）が発生する．

変換器は，以下の手順でSQN管理を実装する．

\begin{enumerate}
\item 5G認証リクエストからSQN⊕AKを抽出（AUTNフィールドの先頭6バイト）
\item 抽出したSQN⊕AKをUEコンテキストにキャッシュ
\item KASME導出時にキャッシュした値を使用
\item 認証失敗時のAUTS（Authentication Token）を処理してSQN再同期
\end{enumerate}

この実装により，認証失敗回数を削減し，アタッチ手順の遅延を最小化できる．

\subsubsection{NAS暗号化と完全性保護}
\paragraph{暗号化アルゴリズムの選択とマッピング}
変換器は，5GのNEA（NAS Encryption Algorithm）と4GのEEA（EPS Encryption Algorithm）間のマッピングを実装する．本研究では，両世代で共通のAES\footnote{Advanced Encryption Standard: 米国標準の共通鍵暗号方式．}暗号プリミティブを使用するNEA2/EEA2を優先的に選択する．

\begin{table}[h]
\centering
\caption{暗号化アルゴリズムのマッピング}
\label{tab:algorithm_mapping}
\begin{tabular}{|l|c|l|c|l|}
\hline
\textbf{5G} & \textbf{ID} & \textbf{4G} & \textbf{ID} & \textbf{暗号プリミティブ} \\
\hline
NEA0 (null) & 0 & EEA0 (null) & 0 & - \\
128-NEA1 & 1 & 128-EEA1 & 1 & SNOW 3G (ストリーム) \\
\textbf{128-NEA2} & \textbf{2} & \textbf{128-EEA2} & \textbf{2} & \textbf{AES-CTR} \\
128-NEA3 & 3 & 128-EEA3 & 3 & ZUC (ストリーム) \\
\hline
\end{tabular}
\end{table}

完全性保護についても同様に，NIA2↔EIA2のマッピングを実装する．両アルゴリズムともAES-CMACを使用するため，無損失変換が可能である．

\paragraph{デュアルモードセキュリティコンテキスト}
変換器は，5G側と4G側で異なる暗号鍵を並列に管理する必要がある．このため，UEごとに以下のセキュリティコンテキストを保持する．

\begin{itemize}
\item \textbf{5G鍵}（AMF通信用）: KNASenc(5G), KNASint(5G), NAS UL COUNT\footnote{NASメッセージの送信回数を示すカウンタ．リプレイ攻撃防止に使用される．}(5G)
\item \textbf{4G鍵}（eNB通信用）: KNASenc(4G), KNASint(4G), NAS UL/DL COUNT(4G)
\item \textbf{アルゴリズム選択}: 5G側（\texttt{selected\_alg\_5g}），4G側（\texttt{selected\_nas\_security\_alg}）
\item \textbf{状態フラグ}: has\_5g\_nas\_keys, has\_nas\_keys
\end{itemize}

この並列鍵管理により，変換器は5G側からの暗号化メッセージを復号し，4G側用に再暗号化する処理を行う．

\paragraph{双方向メッセージ暗号化処理}
セキュリティモード確立後，NASメッセージは暗号化されて転送される．変換器は以下の手順で双方向変換を実行する．

\textbf{ダウンリンク（5G→4G）}:
\begin{enumerate}
\item KNASenc(5G)を使用して5G NASメッセージを復号化（暗号化されている場合）
\item 平文5GMM → EMMメッセージタイプ変換
\item KNASenc(4G)を使用して4G NASメッセージを暗号化（EEA2選択時）
\item KNASint(4G)を使用して完全性保護MAC-I\footnote{Message Authentication Code for Integrity: NASメッセージの改ざん検出に使用される認証コード．}を計算・付加
\end{enumerate}

\textbf{アップリンク（4G→5G）}:
\begin{enumerate}
\item KNASint(4G)を使用して4G NASメッセージの完全性を検証
\item KNASenc(4G)を使用して4G NASメッセージを復号化（暗号化されている場合）
\item 平文EMM → 5GMMメッセージタイプ変換
\item KNASint(5G)を使用して完全性保護MAC-Iを計算・付加
\item （必要に応じて）KNASenc(5G)を使用して5G NASメッセージを暗号化
\end{enumerate}

なお，Initial Context Setup Request等に埋め込まれるNAS-PDUについては，保護ヘッダ付きNASをそのまま保持して転送する方式を採用する．

\subsection{コンテキスト管理とタイムアウト}
\label{subsec:context_management}
変換器はUEごとに以下を保持する．
\begin{itemize}
  \item ID束: (ENB-UE-S1AP-ID, MME-UE-S1AP-ID) ↔ (RAN-UE-NGAP-ID, AMF-UE-NGAP-ID)
  \item セッション情報: PDU Session ID，E-RAB ID，eNB/UPFのトンネル情報（TEID，IPアドレス）
  \item タイムスタンプ: 作成時刻，最終更新時刻，ICS送信時刻等
\end{itemize}
リソース管理として，一定時間未使用のTEIDマッピングやUEコンテキストは定期的にクリーンアップし，メモリリークを防止する．また，手順の重複実行を防ぐためのガード時間を設ける．

異常系（IE不足/不整合等）については，原則としてログ出力の上でメッセージを破棄する．ただし，S1AP/NGAPのErrorIndicationやUE Context Release Request等のエラー関連メッセージを受信した場合は，対向側インターフェースへ変換・転送することで，両端ノード間のエラー通知を橋渡しする．

\subsection{TAU（Tracking Area Update）処理}
\label{subsec:tau_design}
UEが同一TAC内で移動した場合や，定期タイマー（T3412）の満了時に送信されるTAU Requestについては，AMFへの転送ではなく，ゲートウェイがローカルでTAU Acceptを生成して応答する方式を採用する．

この設計の根拠は以下のとおりである．
\begin{itemize}
\item 5GにはTAUに直接対応する手順が存在せず，Registration Updateとして処理される
\item Registration Updateへの変換は，状態管理の複雑化を招く
\item ローカル応答により，コア網への不要なシグナリングを削減できる
\end{itemize}

TAU Acceptには，初期接続時にキャッシュしたTAC/PLMNから生成したTAI List，定期TAUタイマー値，およびアクティブベアラ状態を含める．具体的な実装については\ref{sec:prototype_implementation}節を参照されたい．

\subsection{QoSパラメータのマッピング}
\label{subsec:qos_mapping}
4GのQCI（QoS Class Identifier，TS 23.203）と5GのQoS情報（5QI，TS 23.501）は，数値の定義範囲や付随パラメータが完全には一致しない．

表\ref{tab:qos_mapping}にQCI/5QIの対応例を示す．変換器は，標準仕様においてQCIと5QIが1対1対応する範囲（1〜9）では同一値として扱い，それ以外はデフォルト値へフォールバックするロジックを採用する．

\begin{table}[h]
\centering
\caption{QCIと5QIのマッピング例}
\label{tab:qos_mapping}
\begin{tabular}{|c|c|l|l|}
\hline
\textbf{QCI} & \textbf{5QI} & \textbf{リソースタイプ} & \textbf{用途例} \\
\hline
1 & 1 & GBR & 音声（VoLTE/VoNR） \\
5 & 5 & Non-GBR & IMS Signaling \\
6 & 6 & Non-GBR & ビデオ（バッファ型） \\
9 & 9 & Non-GBR & デフォルト（ベストエフォート） \\
\hline
\end{tabular}
\end{table}

一方で，S1APメッセージ生成時（Downlink，Initial Context Setup Request等）におけるQoS変換については，4G側からのBearer情報と5G側からのQoS Flow情報を適切にマッピングする必要がある．5G側から通知されるQoS Flow ID（QFI）については，コンテキストとして保持し応答メッセージ（PDU Session Resource Setup Response等）に含めることで，プロトコル上の整合性を維持する．

\subsection{エラー処理とロギング}
変換不可能なメッセージや未サポート機能に遭遇した場合，変換器はシステム全体の整合性を維持するため，フェイルセーフ（Fail-safe）な挙動をとる．具体的には以下の処理を行う．

\begin{itemize}
\item \textbf{IE不足/デコード失敗}: 必須パラメータが欠落またはデコードに失敗した場合，ログ出力の上で当該変換を中止する．ただし，NAS変換に失敗した場合は，原文メッセージをそのまま透過転送するフォールバック処理を行う場合もある．
\item \textbf{リソース解放}: 第\ref{subsec:context_management}項で述べた通り，長時間未使用のTEIDマッピングは定期クリーンアップにより解放し，メモリリークを防ぐ．
\item \textbf{障害解析用ロギング}: エラー発生時は，関連ID（ENB-UE-S1AP-ID, RAN-UE-NGAP-ID等）やNAS/手順の種別をログに残し，調査の手がかりとする．
\end{itemize}

\subsection{性能上の考慮事項}
変換器の処理遅延は，主に以下の要素から構成される．
\begin{itemize}
\item \textbf{ASN.1デコード/エンコード}: S1APおよびNGAPメッセージのAPER（Aligned Packed Encoding Rules）処理．一般に計算コストが高い傾向がある．
\item \textbf{NASメッセージ変換}: EMM↔5GMM，ESM↔5GSMのメッセージタイプと情報要素の再構成．
\item \textbf{暗号化処理}: NAS暗号化（NEA/EEA）と完全性保護（NIA/EIA）の計算．ただし，暗号化は常時ではなく条件付き（アルゴリズム選択に依存）で適用される．
\item \textbf{GTP-U TEIDフィールド書き換え}: TEIDマッピングテーブル検索とTEIDフィールド（4バイト）の書き換え．
\item \textbf{コンテキスト検索}: UE IDによるテーブル検索．
\end{itemize}

これらの遅延を低減するため，以下の設計指針が有効である．
\begin{itemize}
\item \textbf{効率的なコンテキスト管理}: UEコンテキストのデータ構造は，メモリ局所性と検索効率を考慮して選択する．
\item \textbf{TEIDマッピングの期限管理}: 第\ref{subsec:context_management}項のリソース管理方針に基づき，長時間未使用のエントリを定期的にクリーンアップする．
\item \textbf{OSネットワークスタックの活用}: GTP-Uパケットの送信にはOSのソケットAPIを使用し，IP/UDPヘッダ生成やチェックサム計算をOSに委ねることで実装を簡素化する．
\end{itemize}

\section{プロトタイプ実装}
\label{sec:prototype_implementation}
前節までに述べた設計を検証するため，本研究ではプロトタイプ実装を行った．本節では，実装の概要と，設計の検証を目的とした本プロトタイプにおける具体的な実装上の選択について述べる．

\subsection{実装概要}
本プロトタイプはC言語で実装し，約14,000行のソースコードから構成される．主要な外部ライブラリとして，S1AP/NGAPメッセージのASN.1エンコード・デコードにはlibasn1c（asn1c生成コード）を使用した．認証処理のMilenageアルゴリズムおよび鍵導出関数（KDF）の実装は，3GPP TS 35.206\cite{3GPP_TS35206}およびTS 33.501の仕様に基づき，Open5GS\cite{open5gs}の実装を参考にした．

表\ref{tab:module_structure}に主要モジュールの構成を示す．

\begin{table}[htbp]
\centering
\caption{プロトタイプの主要モジュール構成}
\label{tab:module_structure}
\begin{tabular}{|l|l|}
\hline
\textbf{モジュール} & \textbf{機能} \\
\hline
core/ & メインループ，SCTP/UDP受信処理 \\
\hline
ngap/ & NGAPメッセージ生成・解析 \\
\hline
nas/ & NASメッセージ変換（4G↔5G） \\
\hline
auth/ & Milenage認証，鍵導出 \\
\hline
context/ & UEコンテキスト・TEIDマッピング管理 \\
\hline
transport/ & GTP-Uパケット転送 \\
\hline
\end{tabular}
\end{table}

以下では，本プロトタイプにおける具体的な実装上の選択と簡略化について述べる．

\subsection{ID変換における簡略化}
gNB-ID生成については，eNB-IDからの動的変換ではなく，固定値を用いて生成している．同様に，TACについてもS1SetupRequestから取得できない場合は環境変数により設定したデフォルト値へフォールバックする構成とした．

\subsection{Early S1 Setupにおける固定値}
PDUセッション確立前の仮TEID割り当てには，固定値\texttt{0x00000001}を使用している．また，PDU Session Establishment Requestで使用するAPNについても，UEから受信したAttach RequestのESM情報の解析は行わず，環境変数で指定されたデフォルト値（\texttt{internet}）を使用する構成とした．

\subsection{U-PlaneのTransparent Proxy動作}
本実装では，Downlink方向においてTransparent proxy方式を採用し，TEID変換を行っていない．UPFがeNBのS1-U TEIDを直接使用して送出するため，ゲートウェイはTEIDを書き換えず宛先IPアドレスのみをeNBに変更する．

\subsection{認証・鍵管理における構成}
検証環境におけるKi/OPcの取得には，設定ファイル（auth\_keys.yaml）からの読み込み方式を採用した．この構成に伴うセキュリティ上の制約については\ref{subsec:security_constraints}節で述べる．

\subsection{暗号化アルゴリズムの対応範囲}
本実装ではNEA2/EEA2（AES系）を中心に対応しており，NEA1/EEA1（SNOW 3G）およびNEA3/EEA3（ZUC）は未実装である．NAS暗号化は双方向（Uplink/Downlink）で実装されており，4G側・5G側それぞれの鍵を用いた暗号化・復号処理を行う．

\subsection{TAU処理における実装}
\ref{subsec:tau_design}節で述べたTAU処理のローカル応答方式を実装した．TAU Acceptの生成には，初期接続時にキャッシュしたTAC/PLMN情報を使用する．

\subsection{タイムアウト値の設定}
本実装では，汎用的なタイマ機構ではなくタイムスタンプベースの時間管理を採用し，以下の固定値を設定している．
\begin{itemize}
  \item \textbf{TEIDマッピングの期限切れ}: 1時間未使用で解放（5分間隔のクリーンアップ）
  \item \textbf{TAUガード}: 60秒のガード時間でフラグ管理（連続TAU抑止）
  \item \textbf{ICS重複抑止}: 10秒間の重複送信抑止
\end{itemize}

\subsection{QoSマッピングの簡略化}
S1APメッセージ生成時（Downlink方向）においては，接続性の検証を主目的とするため動的なQoS変換は行わず，一律でデフォルト値（QCI=9）を固定で設定する構成とした．

\subsection{データ構造の選択}
UEコンテキストは固定長配列（\texttt{ue\_mappings}）で管理し，線形探索により検索する．実装は簡素であり，配列のメモリ局所性によりキャッシュ効率が期待できるが，UE数が増加した場合のスケーラビリティには制約がある．

\subsection{論理設計との差分}
表\ref{tab:logical_vs_prototype}に，第\ref{chap:proposal}章の論理設計と本プロトタイプ実装の主な差分を整理する．

\begin{table}[htbp]
\centering
\caption{論理設計とプロトタイプ実装の差分}
\label{tab:logical_vs_prototype}
\begin{tabular}{|l|p{5cm}|p{5cm}|}
\hline
\textbf{項目} & \textbf{論理設計} & \textbf{プロトタイプ実装} \\
\hline
gNB-ID生成 & eNB-IDからの動的変換 & 固定値による生成 \\
\hline
APN取得 & ESM情報の解析により取得 & 環境変数によるデフォルト値 \\
\hline
TEID変換 & 双方向での動的マッピング & Downlink方向はTransparent proxy \\
\hline
Ki/OPc取得 & セキュアな外部連携 & 設定ファイルからの読み込み \\
\hline
暗号アルゴリズム & 全アルゴリズム対応 & NEA2/EEA2（AES系）のみ \\
\hline
QoSマッピング & QCI/5QIの動的変換 & QCI=9固定 \\
\hline
UEコンテキスト管理 & スケーラブルなデータ構造 & 固定長配列・線形探索 \\
\hline
\end{tabular}
\end{table}

これらの差分は，本プロトタイプが接続性検証を主目的としたPoC（Proof of Concept）であることに起因する．商用実装においては，各項目について論理設計に沿った実装が必要となる．

\section{プロトコル仕様上の制約}
\label{sec:protocol_constraints}
本節では，S1AP/NGAP間のプロトコル変換において生じる本質的な制約を整理する．

\subsection{規格差分による制約}
S1APとNGAPは情報要素（IE）の集合と手順が一致しないため，完全な1対1変換は成立しない．S1APに対応する要素が存在しないIE（例: AMF Region情報）は，S1AP側へは反映されない．また，NGAPで必須となるNetwork Slice関連IE（AllowedNSSAI，S-NSSAI等）については，4G側に対応概念がないため，固定値（SST=1等）での補完が必要となる（\ref{sec:prototype_implementation}節参照）．

必須IEの欠落など致命的な不整合では当該変換を中止するが，状況により原文メッセージを透過転送するフォールバックも用いる．

\subsection{手順分離に起因する制約}
5GではRegistration完了後にPDU Session確立が行われる一方，4GではAttach手順の一部としてデータ経路（Default Bearer）が確立される．この差分に対して，本研究では\S\ref{sec:pdu_session_establishment}で述べたEarly S1 Setup手順により，手順の逐次化と待ち合わせを行う．

ただし，PDU Session確立に必要な情報（TEID等）が通知されるまでの待機が発生する．コア網の処理遅延が大きい環境では，確立成功率の低下や再試行による遅延増加が生じる可能性がある．

\subsection{機能制限に関する制約}
本変換器設計は以下の機能について本質的な制約を持つ．これらは第\ref{chap:proposal}章の論理設計で対応が求められるとした機能のうち，本プロトタイプでは実装対象外としたもの，または簡易実装に留めたものである．
\begin{itemize}
\item \textbf{ハンドオーバー}: S1ハンドオーバー手順とN2ハンドオーバー手順の相互変換は未実装であり，基地局間の移動を伴う通信継続には対応していない．本研究では位置登録（Registration）に焦点を当て，移動管理（Mobility Management）は対象外としている．
\item \textbf{Network Slicing}: S-NSSAIに基づくスライス選択は4Gに対応概念がないため，NGAP必須IEを満たすために固定値（SST=1等）での補完が必要となる．論理設計ではAPNからS-NSSAIへのマッピングを想定しているが，本PoCでは未実装である．
\item \textbf{QoS Flow}: 5GのFlow単位QoS\allowbreak （GFBR/MFBR）から4GのBearer単位\allowbreak （GBR/MBR）への完全なマッピングは困難であり，QCI/5QI/QFIの対応関係に基づく簡易変換（\ref{subsec:qos_mapping}節参照）に留まる．
\item \textbf{Traffic Steering}: ULCL（Uplink Classifier）等による動的なトラフィック分岐は，4GにULCLに対応する概念がなく，本PoCでは対象外としている．
\item \textbf{Multi-type PDU Support}: 5GCはIPに加えEthernetやUnstructuredタイプをサポートするが，4G RANはIPベースのPDN接続のみであり，本PoCではIPv4/IPv6のみを対象としている．
\end{itemize}

\subsection{セキュリティに関する制約}
\label{subsec:security_constraints}
\ref{subsec:auth_key_management}節で述べたとおり，本提案の変換器設計ではUEの秘密情報（Ki, OPc）へのアクセス手段が必要となる．N26インターフェースを介した標準的な鍵共有が利用できない構成では，変換器が何らかの形でこれらの情報を保持する必要があり，セキュリティアーキテクチャ上の考慮が必要となる．

本研究の検証環境では，設定ファイル（auth\_keys.yaml）からKi/OPcを読み込む方式を採用した．これは商用環境では一般に想定されない構成であり，実運用においては以下のいずれかのアプローチが考えられる．
\begin{itemize}
\item 標準のN26インターフェースを実装し，コア網間で鍵コンテキストを共有する
\item HSS/UDM連携により，変換器が認証ベクトルを取得する
\item 変換器を信頼されたセキュリティドメイン内に配置し，適切なアクセス制御を実施する
\end{itemize}
\section{本章のまとめ}

本章では，第\ref{chap:proposal}章で提案した疎結合アーキテクチャを実現するための詳細設計，プロトタイプ実装，および制約事項を述べた．各節の貢献を以下に要約する．

\begin{itemize}
    \item \textbf{詳細設計}（第\ref{sec:detailed_design}節）: C-Plane変換（S1AP↔NGAP，NAS EPS↔NAS 5GSのメッセージ変換規則，ID変換の体系化），U-Plane変換（TEIDマッピングテーブルによるHybrid Proxy設計），認証・鍵管理（Milenage処理代行によるデュアルモードセキュリティ），および手順タイミング差異を解消するEarly S1 Setup手順を設計した．
    \item \textbf{プロトタイプ実装}（第\ref{sec:prototype_implementation}節）: C言語で約14,000行のプロトタイプを実装し，接続性検証を主目的としたPoCとしての簡略化（gNB-ID固定値生成，APN環境変数設定，QCI=9固定等）と論理設計との差分を明確化した．
    \item \textbf{プロトコル仕様上の制約}（第\ref{sec:protocol_constraints}節）: 規格差分による制約（S1AP/NGAP間のIE非対称性），手順分離に起因する制約（PDU Session確立待機），機能制限に関する制約（ハンドオーバー，Network Slicing，QoS Flow，Traffic Steering，Multi-type PDU），およびセキュリティに関する制約（Ki/OPcアクセス要件）を整理した．
\end{itemize}

プロトタイプ実装はC言語で約14,000行，概念実証を目的とした構成である．次章では，このプロトタイプを用いた実験環境と評価手法を述べる．
